\section{Задание}
Для данной таблицы значений функции построить методом наименьших квадратов линейную и квадратичную полиномиальные модели, вычислить наилучшие среднеквадратические отклонения.

\section{Цель работы}
Научиться строить полиномиальные аналитические модели табличных функций методом наименьших квадратов.

\section{Ход работы}
\textit{Перед выполнением работы отметим, что все исходные данные взяты из варианта 47.} Исходные данные представлены в таблице \ref{table:values}.

\begin{table}[H]
    \centering
    \caption{Исходные значения $x_i, y_i$ (вариант 47)}
    \label{table:values}
    \begin{tblr}{
        colspec = {X[-1,l]X[c]X[c]X[c]X[c]X[c]X[c]X[c]X[c]X[c]X[c]},
        hlines,
        vlines,
        rows = {abovesep=5pt, belowsep=5pt, font=\small, $},
        column{1} = {leftsep=8pt},
    }
        x_i & -2.0 & -1.9 & -1.8 & -1.7 & -1.6 & -1.5 & -1.4 & -1.3 & -1.2 & -1.1 \\
        y_i & 12.28 & 12.53 & 12.50 & 12.53 & 12.75 & 12.85 & 12.77 & 12.76 & 12.73 & 12.85 \\
    \end{tblr}

    \vspace{6pt}

    \begin{tblr}{
        colspec = {X[-1,l]X[c]X[c]X[c]X[c]X[c]X[c]X[c]X[c]X[c]X[c]},
        hlines,
        vlines,
        rows = {abovesep=5pt, belowsep=5pt, font=\small, $},
        column{1} = {leftsep=8pt},
    }
        x_i & -1.0 & -0.9 & -0.8 & -0.7 & -0.6 & -0.5 & -0.4 & -0.3 & -0.2 & -0.1 \\
        y_i & 12.51 & 12.34 & 12.22 & 11.84 & 11.67 & 11.27 & 11.06 & 10.73 & 10.35 & 10.09 \\
    \end{tblr}
\end{table}

\subsection{Метод наименьших квадратов}

Пусть функция $y = f(x)$ задана таблицей значений $\{(x_i, y_i)\}_{i=1}^{n}$. Требуется найти полином $P_m(x) = a_0 + a_1 x + \dots + a_m x^m$, наилучшим образом приближающий табличные данные. Коэффициенты $a_0, a_1, \dots, a_m$ определяются из условия минимума суммы квадратов отклонений:
\[
S(a_0, \dots, a_m) = \sum_{i=1}^{n} \left(y_i - \sum_{j=0}^{m} a_j x_i^j \right)^2 \to \min.
\]

Приравнивая частные производные $\frac{\partial S}{\partial a_k}$ к нулю, получаем нормальную систему:
\[
\sum_{j=0}^{m} a_j \sum_{i=1}^{n} x_i^{j+k} = \sum_{i=1}^{n} y_i x_i^k, \quad k = 0, \dots, m.
\]

Наилучшее среднеквадратическое отклонение вычисляется по формуле
\[
\delta_2(P_m) = \sqrt{\frac{1}{n} \sum_{i=1}^{n} (y_i - P_m(x_i))^2}.
\]

\subsection{Линейная аппроксимация}

Ищем $P_1(x) = a_0 + a_1 x$. Нормальная система принимает вид
\[
\begin{cases}
a_0 n + a_1 \displaystyle\sum_{i=1}^{n} x_i = \displaystyle\sum_{i=1}^{n} y_i, \\[10pt]
a_0 \displaystyle\sum_{i=1}^{n} x_i + a_1 \displaystyle\sum_{i=1}^{n} x_i^2 = \displaystyle\sum_{i=1}^{n} y_i \, x_i.
\end{cases}
\]

\noindentРешаем систему методом Крамера. Определитель матрицы системы:
\[
\Delta = \begin{vmatrix}
n & \sum x_i \\
\sum x_i & \sum x_i^2
\end{vmatrix}.
\]

\noindentОпределители для нахождения коэффициентов:
\[
\Delta_0 = \begin{vmatrix}
\sum y_i & \sum x_i \\
\sum y_i x_i & \sum x_i^2
\end{vmatrix}, \quad
\Delta_1 = \begin{vmatrix}
n & \sum y_i \\
\sum x_i & \sum y_i x_i
\end{vmatrix}.
\]

\noindentКоэффициенты:
\[
a_0 = \frac{\Delta_0}{\Delta}, \quad a_1 = \frac{\Delta_1}{\Delta}.
\]

\noindentНа листинге \ref{lst:linear} приведена программа, строящая линейную модель МНК.

\vspace{10pt}
\captionof{listing}{Построение линейной модели МНК}
\label{lst:linear}
\vspace{-10pt}
\begin{minted}[linenos, numbersep=10pt, firstnumber=1]{matlab}
x = [-2.0, -1.9, -1.8, -1.7, -1.6, -1.5, -1.4, -1.3, -1.2, -1.1, ...
     -1.0, -0.9, -0.8, -0.7, -0.6, -0.5, -0.4, -0.3, -0.2, -0.1];
y = [12.28, 12.53, 12.50, 12.53, 12.75, 12.85, 12.77, 12.76, 12.73, 12.85, ...
     12.51, 12.34, 12.22, 11.84, 11.67, 11.27, 11.06, 10.73, 10.35, 10.09];

n = length(x);

A = [n,        sum(x);
     sum(x),   sum(x.^2)];
b = [sum(y); sum(x .* y)];

D = det(A);

A0 = A; A0(:, 1) = b;
A1 = A; A1(:, 2) = b;

a0 = det(A0) / D;
a1 = det(A1) / D;

fprintf('P_1(x) = %.4f + (%.4f) * x\n', a0, a1);

delta = sqrt(1/n * sum((y - a0 - a1 * x).^2));
fprintf('delta_2(P_1) = %.4f\n', delta);
\end{minted}

\noindentНа рис. \ref{fig:task1a.png} предаствлен результат работы программы.

\screenshot[.6]{task1a.png}{Результат работы программы, строящей линейную модель МНК}

\noindentВ результате работы программы были получкны следующие коэффициенты:

\[
a_0 = 10.7529, \quad a_1 = - 1.2177.
\]

\noindentЛинейная модель выглядит так:

\[
P_1(x) = 10.75 - 1.22x.
\]

\noindentНаилучшее среднеквадратическое отклонение:
\[
\delta_2(P_1) = 0.48.
\]

\subsection{Квадратичная аппроксимация}

Ищем $P_2(x) = a_0 + a_1 x + a_2 x^2$. Нормальная система:
\[
\begin{cases}
a_0 n + a_1 \displaystyle\sum_{i=1}^{n} x_i + a_2 \displaystyle\sum_{i=1}^{n} x_i^2 = \displaystyle\sum_{i=1}^{n} y_i, \\[10pt]
a_0 \displaystyle\sum_{i=1}^{n} x_i + a_1 \displaystyle\sum_{i=1}^{n} x_i^2 + a_2 \displaystyle\sum_{i=1}^{n} x_i^3 = \displaystyle\sum_{i=1}^{n} y_i \, x_i, \\[10pt]
a_0 \displaystyle\sum_{i=1}^{n} x_i^2 + a_1 \displaystyle\sum_{i=1}^{n} x_i^3 + a_2 \displaystyle\sum_{i=1}^{n} x_i^4 = \displaystyle\sum_{i=1}^{n} y_i \, x_i^2.
\end{cases}
\]

\noindentРешаем систему методом Крамера. Определитель матрицы системы:
\[
\Delta = \begin{vmatrix}
n & \sum x_i & \sum x_i^2 \\
\sum x_i & \sum x_i^2 & \sum x_i^3 \\
\sum x_i^2 & \sum x_i^3 & \sum x_i^4
\end{vmatrix}.
\]

\noindentОпределители для нахождения коэффициентов:
\begin{gather*}
\Delta_0 = \begin{vmatrix}
\sum y_i & \sum x_i & \sum x_i^2 \\
\sum y_i x_i & \sum x_i^2 & \sum x_i^3 \\
\sum y_i x_i^2 & \sum x_i^3 & \sum x_i^4
\end{vmatrix}, \quad
\Delta_1 = \begin{vmatrix}
n & \sum y_i & \sum x_i^2 \\
\sum x_i & \sum y_i x_i & \sum x_i^3 \\
\sum x_i^2 & \sum y_i x_i^2 & \sum x_i^4
\end{vmatrix}, \\[5pt]
\Delta_2 = \begin{vmatrix}
n & \sum x_i & \sum y_i \\
\sum x_i & \sum x_i^2 & \sum y_i x_i \\
\sum x_i^2 & \sum x_i^3 & \sum y_i x_i^2
\end{vmatrix}.
\end{gather*}

Коэффициенты:
\[
a_0 = \frac{\Delta_0}{\Delta}, \quad a_1 = \frac{\Delta_1}{\Delta}, \quad a_2 = \frac{\Delta_2}{\Delta}.
\]

На листинге \ref{lst:quadratic} приведена программа, строящая квадратичную модель МНК.

\vspace{10pt}
\captionof{listing}{Построение квадратичной модели МНК}
\label{lst:quadratic}
\vspace{-10pt}
\begin{minted}[linenos, numbersep=10pt, firstnumber=1]{matlab}
x = [-2.0, -1.9, -1.8, -1.7, -1.6, -1.5, -1.4, -1.3, -1.2, -1.1, ...
     -1.0, -0.9, -0.8, -0.7, -0.6, -0.5, -0.4, -0.3, -0.2, -0.1];
y = [12.28, 12.53, 12.50, 12.53, 12.75, 12.85, 12.77, 12.76, 12.73, 12.85, ...
     12.51, 12.34, 12.22, 11.84, 11.67, 11.27, 11.06, 10.73, 10.35, 10.09];

n = length(x);

A = [n,         sum(x),    sum(x.^2);
     sum(x),    sum(x.^2), sum(x.^3);
     sum(x.^2), sum(x.^3), sum(x.^4)];
b = [sum(y); sum(x .* y); sum(x.^2 .* y)];

D = det(A);

A0 = A; A0(:, 1) = b;
A1 = A; A1(:, 2) = b;
A2 = A; A2(:, 3) = b;

a0 = det(A0) / D;
a1 = det(A1) / D;
a2 = det(A2) / D;

fprintf('P_2(x) = %.4f + (%.4f) * x + (%.4f) * x^2\n', a0, a1, a2);

delta = sqrt(1/n * sum((y - a0 - a1 * x - a2 * x.^2).^2));
fprintf('delta_2(P_2) = %.4f\n', delta);
\end{minted}

\noindentНа рис. \ref{fig:task1b.png} предаствлен результат работы программы.

\screenshot[.8]{task1b.png}{Результат работы программы, строящей квадратичную модель МНК}

\noindentВ результате работы программы были получкны следующие коэффициенты:

\[
a_0 = 9.5155, \quad a_1 = -4.5924, \quad a_2 = -1.6070.
\]


\noindentКвадратичная модель выглядит так:
\[
P_2(x) = 9.52 - 4.59x - 1.61x^2.
\]

\noindentНаилучшее среднеквадратическое отклонение:
\[
\delta_2(P_2) = 0.09.
\]

\section{Вывод}

В результате выполнения лабораторной работы были построены линейная и квадратичная полиномиальные модели табличной функции методом наименьших квадратов. Для каждой модели были найдены коэффициенты путём решения соответствующей нормальной системы уравнений методом Крамера и вычислены наилучшие среднеквадратические отклонения: $\delta_2(P_1) = 0.48$ для линейной модели и $\delta_2(P_2) = 0.09$ для квадратичной. Квадратичная модель даёт значительно меньшее отклонение, что говорит о лучшем приближении табличных данных.
